Thinking...
The user asks: "以下のの日本語文を英訳してください。ただしTeXのコードが挿入されているのでその部分はそのまま出力してください：" They give a Japanese paragraph with LaTeX references: 表\ref{tab:G1topj}-\ref{tab:G4topj}, Journal Citation Reports (JCR）\cite{JCR1}, 図\ref{fig:gr12}. Need to translate to English, preserving TeX code unchanged. Ensure proper punctuation. Let's produce translation.

We must keep the LaTeX commands as is: \ref{tab:G1topj}, \ref{tab:G4topj}, \cite{JCR1}, \ref{fig:gr12}. Also maybe parentheses.

Let's produce English translation:

"Here we compare the changes in journals that publish highly cited papers with those of the non‑highly‑cited group. Tables \ref{tab:G1topj}–\ref{tab:G4topj} list the journals in which the top‑ranked papers of groups G1–G4 appear. The impact factor of each journal and its year of publication were obtained from the Journal Citation Reports (JCR) \cite{JCR1}. Because JCR has published impact factors only since 1997, we calculated the impact factor for the final year of each cumulative cohort only for cohorts from 1946–2000 onward. From this we see that journals that publish highly cited papers are not necessarily journals with high impact factors. On the other hand, there is no special difference in journal composition between groups. Therefore, to examine the degree of journal preference within each group, we computed the Kullback–Leibler divergence (KLD) between the combined population of all groups and each individual group, as shown in Figure \ref{fig:gr12}. This figure reveals that the highly cited papers exhibit a stronger journal preference than the other groups."

Make sure TeX code remains unchanged. Provide just the translation.
...done thinking.

Here, we compare the changes in journals that publish highly cited papers with those of the non‑highly‑cited group. Tables \ref{tab:G1topj}–\ref{tab:G4topj} show the list of journals in which the top‑ranked papers of groups G1–G4 appear. The impact factor of each journal and its year of publication were obtained from the Journal Citation Reports (JCR) \cite{JCR1}. Since JCR has released impact factors only since 1997, we calculated the impact factor for the final year of each cumulative cohort only for the cohorts from 1946–2000 onward. From this, we see that journals that publish highly cited papers are not necessarily journals with high impact factors. On the other hand, no special differences in journal composition are observed between the groups. Therefore, to assess the degree of journal preference in each group, we computed the Kullback‑Leibler divergence (KLD) between the pooled population of all groups and each individual group, as shown in Figure \ref{fig:gr12}. This figure makes it clear that highly cited papers have a stronger journal preference compared with the other groups.

